\clearpage
\chapter*{Fe de erratas de la exposición n\textsuperscript{o} 182}
\addcontentsline{toc}{chapter}{Fe de erratas de la exposición n\textsuperscript{o} 182}
\setcounter{footnote}{1}
\src{298}

\noindent\textbf{2. Geometría formal y geometría algebraica.}

\begin{center}
[Seminario Bourbaki, t. 11, 1958/59, n\textsuperscript{o} 182, 28 p.]
\end{center}

\noindent\textbf{Página 8, teorema 6.} --- La hipótesis «$f$
cuasiproyectivo» puede sustituirse por la hipótesis más débil «$f$
separado», en virtud del resultado siguiente (cf. SGA VIII, 6.2): todo morfismo
$f:X\longrightarrow Y$ que es cuasifinito y separado es cuasiproyectivo.

\src{299}
\noindent\textbf{Página 12, y página 14, observación 1.} --- Señalemos que J.-P.
Serre ha construido una variedad proyectiva no singular, de dimensión 3, sobre
un cuerpo algebraicamente cerrado $k$, de característica $p>0$, que no proviene
por reducción de un esquema propio sobre un anillo local íntegro de cuerpo
residual $k$ y con cuerpo de fracciones de característica
cero.\footnote{Serre (Jean-Pierre). --- Exemples de variétés projectives en
caractéristique $p$ non relevables en caractéristique zéro, \emph{Proc. Nat.
Acad. Sc. U.S.A.}, t. 47, 1961, p. 108--109.} Mumford habría hallado un ejemplo
análogo, con una superficie proyectiva no singular.

\noindent\textbf{Página 15, observaciones 2 y 3.} --- Actualmente soy menos
optimista respecto de los resultados aquí conjeturados. Sin embargo, el
problema relativo a las variaciones de estructura del espacio proyectivo,
señalado al final de la observación 3, ha sido resuelto afirmativamente por
Hironaka, y el problema análogo para las variedades abelianas ha sido resuelto
por Koizumi.

\noindent\textbf{Página 16, línea 3.} --- Señalemos que Mumford ha construido
recientemente los esquemas de módulos para las curvas de género $g$ (cf. SHMI).
El teorema 10 prueba, por lo demás, que los «esquemas de Jacobi de nivel $n$»
de la teoría de Módulos son no singulares (e incluso simples sobre $\mathbf Z$).

\noindent\textbf{Página 23, línea 12, primera línea del corolario
5.} --- Añadir: «siendo $X$ o $Y$ propio sobre $k$».

La sustancia de los §§ 1 a 5 está contenida en la parte publicada de EGA III, y
la de los §§ 6 y 7 está contenida en SGA III. Para el estudio del grupo
fundamental, véanse SGA V, IX, X, XI, así como SGA 1962 (exposiciones X, XII y
XIII), para los teoremas de tipo Lefschetz y numerosas cuestiones abiertas.
Solo la teoría de los revestimientos moderadamente ramificados (cf. teorema 14)
no ha sido todavía objeto de una redacción formal. El corolario del teorema 14,
que da la determinación completa de los revestimientos galoisianos de orden
primo con la característica de una curva algebraica sobre un cuerpo
algebraicamente cerrado, ha sido utilizado de manera esencial en tres ocasiones:

1\textsuperscript{o} en la demostración por Igusa de la desigualdad de Picard
para las superficies proyectivas no singulares en característica cualquiera;

2\textsuperscript{o} en el estudio (desarrollado independientemente por Ogg y
Šafarevič) del grupo de los fibrados principales homogéneos bajo una variedad
abeliana definida sobre un cuerpo de funciones de una variable, en
característica cualquiera;

3\textsuperscript{o} en la demostración reciente, por Artin, de ciertos
teoremas clave sobre la «cohomología de Weil» de las variedades algebraicas.
